% ╔══════════════════════════════════════════════════════════════════════╗
% ║  LICEO V.A.L.  –  DEPARTAMENTO DE MATEMÁTICAS                        ║
% ║  Examen Final: Unidad 17 (06M-17 FL) |  Grado: 6°                    ║
% ╚══════════════════════════════════════════════════════════════════════╝
\documentclass[12pt,letterpaper]{article}

% ── Codificación y tipografía ──────────────────────────────────────────
\usepackage{lmodern}
\usepackage[utf8]{inputenc}
\usepackage[T1]{fontenc}
\usepackage[spanish, es-noshorthands]{babel}
\usepackage{helvet}
\renewcommand{\familydefault}{\sfdefault}

% ── Geometría ─────────────────────────────────────────────────────────
\usepackage[letterpaper,left=1.5cm,right=1.5cm,top=1.5cm,bottom=1.5cm]{geometry}

% ── Gráficos, color y símbolos ────────────────────────────────────────
\usepackage[table]{xcolor}
\usepackage{graphicx}
\usepackage{pifont}
\usepackage{tikz}
\usetikzlibrary{shapes.geometric,arrows.meta,calc,positioning}
\usepackage[most]{tcolorbox} % Modificado para evitar error de librería
\usepackage{array,tabularx,multicol,enumitem}

% ── Paleta de Colores Liceo VAL ───────────────────────────────────────
\definecolor{valblue}       {RGB}{  0, 84,166}
\definecolor{valgreen}      {RGB}{  0,150, 64}
\definecolor{valred}        {RGB}{200, 16, 46}
\definecolor{valdark}       {RGB}{ 40, 40, 40}
\definecolor{valorange}     {RGB}{230,100,  0}
\definecolor{valgray}       {RGB}{245,245,245}
\definecolor{vallightblue}  {RGB}{220,235,250} % Agregado para corregir el error

% ── Casilla de respuesta ──────────────────────────────────────────────
\newcommand{\boxfill}[1][1.5cm]{%
	\tikz[baseline=-0.5ex]%
	\draw[valdark!60, fill=valgray, line width=0.8pt, rounded corners=1pt] (0,-0.3) rectangle (#1,0.3cm);%
}

\pagestyle{empty}
\setlength{\parindent}{0pt}

% ============================================================
\begin{document}
	
	% ─────────────────────────────────────────────────────────────
	% ENCABEZADO DEL EXAMEN
	% ─────────────────────────────────────────────────────────────
	\begin{center}
		{\Large\bfseries\color{valblue} LICEO V.A.L.\ (VIDA, AMOR, LUZ)}\\[2pt]
		{\large\bfseries DEPARTAMENTO DE MATEMÁTICAS -- GRADO 6\textdegree}\\[6pt]
		\begin{tcolorbox}[enhanced,arc=3mm,boxrule=1.5pt,
			colback=valblue,colframe=valdark,halign=center, top=4pt, bottom=4pt]
			{\Large\bfseries\color{white} EXAMEN FINAL: UNIDAD 17}\\[2pt]
			{\large\bfseries\color{yellow!90!white} Estadística: Lectura de tablas, gráficas y pictogramas}
		\end{tcolorbox}
		\vspace{2pt}
		\begin{tcolorbox}[enhanced,arc=2mm,boxrule=1pt,colback=white,colframe=valblue,width=\linewidth]
			\textbf{Nombre:} \rule{8cm}{0.5pt} \hfill \textbf{Fecha:} \rule{3cm}{0.5pt} \hfill \textbf{Nota:} \boxfill[1.5cm]
		\end{tcolorbox}
	\end{center}
	\vspace{-4pt}
	
	% ─────────────────────────────────────────────────────────────
	% PUNTO 1: TABLAS
	% ─────────────────────────────────────────────────────────────
	{\large\bfseries\color{valblue} 1. Análisis de Tablas (25\%)}\\[4pt]
	La siguiente tabla muestra la cantidad de material reciclable (en kg) recolectado por tres cursos.
	
	\begin{minipage}[c]{0.45\linewidth}
		\centering
		\renewcommand{\arraystretch}{1.4}
		\begin{tabular}{|l|c|c|c|}
			\hline
			\rowcolor{valblue!20}
			\textbf{Curso} & \textbf{Cartón} & \textbf{Plástico} & \textbf{Vidrio} \\ \hline
			Sexto A & 15 & 8  & 10 \\ \hline
			Sexto B & 12 & 14 & 6  \\ \hline
			Sexto C & 20 & 5  & 9  \\ \hline
		\end{tabular}
	\end{minipage}%
	\hfill
	\begin{minipage}[c]{0.52\linewidth}
		\begin{enumerate}[label=\alph*), itemsep=0.3cm, leftmargin=*]
			\small
			\item ¿Cuántos kg de \textbf{Plástico} recogió Sexto B? \hfill\boxfill
			\item ¿Qué curso recogió más \textbf{Cartón}? \hfill\boxfill[2cm]
			\item Sumando los tres cursos, ¿cuánto \textbf{Vidrio} hay en total? \hfill\boxfill
			\item ¿Cuántos kg recogió \textbf{Sexto A en total} (sumando sus 3 materiales)? \hfill\boxfill
		\end{enumerate}
	\end{minipage}
	
	\vspace{12pt}\noindent\rule{\linewidth}{1pt}\vspace{12pt} % Modificado para evitar error de \hrule
	
	% ─────────────────────────────────────────────────────────────
	% PUNTO 2: BARRAS
	% ─────────────────────────────────────────────────────────────
	{\large\bfseries\color{valorange} 2. Gráfica de Barras (25\%)}\\[4pt]
	Observa los puntos anotados por 4 equipos en un torneo. \textit{(Recuerda usar la línea guía hacia el eje Y)}.
	
	\begin{minipage}[c]{0.48\linewidth}
		\centering
		\begin{tikzpicture}[x=1.2cm,y=0.12cm]
			% Cuadrícula cada 5 unidades
			\foreach \y in {0,5,10,15,20,25,30} {
				\draw[gray!30, line width=0.5pt] (0,\y) -- (5,\y);
				\node[left, font=\bfseries\scriptsize, color=valdark] at (0,\y) {\y};
			}
			\draw[valdark, line width=1.2pt] (0,0) -- (0,32);
			\draw[valdark, line width=1.2pt] (0,0) -- (5.2,0);
			\node[rotate=90, font=\bfseries\scriptsize, color=valdark] at (-0.8, 15) {Puntos Anotados};
			
			% Barras
			\foreach \x/\y/\col/\name in {
				1/20/valorange/Tigres,
				2/15/valorange/Leones,
				3/30/valorange/Pumas,
				4/5/valorange/Osos%
			} {
				\path[fill=\col!75, draw=\col, line width=1pt] (\x-0.3,0) rectangle (\x+0.3,\y);
				\node[below, font=\bfseries\scriptsize, color=valdark] at (\x,0) {\name};
			}
		\end{tikzpicture}
	\end{minipage}%
	\hfill
	\begin{minipage}[c]{0.50\linewidth}
		\begin{enumerate}[label=\alph*), itemsep=0.35cm, leftmargin=*]
			\small
			\item Puntos exactos del equipo \textbf{Leones}: \hfill\boxfill
			\item Equipo que anotó \textbf{30 puntos}: \hfill\boxfill[2cm]
			\item ¿Cuántos puntos \textbf{más} hicieron los Tigres que los Osos? \hfill\boxfill
			\item ¿Cuántos puntos anotaron los 4 equipos \textbf{en total}? \hfill\boxfill
		\end{enumerate}
	\end{minipage}
	
	\vspace{12pt}\noindent\rule{\linewidth}{1pt}\vspace{12pt}
	
	% ─────────────────────────────────────────────────────────────
	% PUNTO 3: LÍNEAS
	% ─────────────────────────────────────────────────────────────
	{\large\bfseries\color{valgreen} 3. Gráfica de Líneas (25\%)}\\[4pt]
	La gráfica muestra la temperatura registrada en una ciudad durante cinco días a las 2:00 PM. 
	
	\begin{minipage}[c]{0.52\linewidth}
		\centering
		\begin{tikzpicture}[x=1.2cm,y=0.25cm]
			% Cuadrícula cada 2 unidades (Evitando cálculos matemáticos en la coordenada)
			\foreach \y/\ylab in {0/10, 2/12, 4/14, 6/16, 8/18, 10/20, 12/22, 14/24} {
				\draw[gray!30, line width=0.5pt] (0,\y) -- (5.5,\y);
				\node[left, font=\bfseries\scriptsize, color=valdark] at (0,\y) {\ylab};
			}
			\draw[valdark, line width=1.2pt] (0,0) -- (0,15);
			\draw[valdark, line width=1.2pt] (0,0) -- (5.7,0);
			
			% Modificado para incluir el símbolo de grados de forma segura
			\node[rotate=90, font=\bfseries\scriptsize, color=valdark] at (-0.8, 7.5) {Temperatura ($^\circ$C)};
			
			% Eje X marcas
			\foreach \x/\name in {1/Lun, 2/Mar, 3/Mié, 4/Jue, 5/Vie} {
				\draw[valdark, line width=1pt] (\x,0) -- (\x,-0.5);
				\node[below, font=\bfseries\scriptsize] at (\x,-0.5) {\name};
			}
			
			% Línea y puntos (Valores escalados: ej. 14 -> 4, 20 -> 10)
			\draw[valgreen, line width=1.5pt] (1, 4) -- (2, 10) -- (3, 6) -- (4, 6) -- (5, 12);
			\foreach \x/\y in {1/4, 2/10, 3/6, 4/6, 5/12} {
				\filldraw[valdark, draw=white, line width=1pt] (\x,\y) circle (2.5pt);
			}
		\end{tikzpicture}
	\end{minipage}%
	\hfill
	\begin{minipage}[c]{0.45\linewidth}
		\begin{enumerate}[label=\alph*), itemsep=0.35cm, leftmargin=*]
			\small
			\item Temperatura el día \textbf{Martes}: \hfill\boxfill
			\item Día con la temperatura \textbf{más alta}: \hfill\boxfill[2cm]
			\item Entre el martes y el miércoles, ¿la temperatura \textbf{subió o bajó}? \hfill\boxfill[2cm]
			\item ¿En qué días la temperatura se mantuvo \textbf{constante} (igual)? \hfill\boxfill[2cm]
		\end{enumerate}
	\end{minipage}
	
	\vspace{12pt}\noindent\rule{\linewidth}{1pt}\vspace{12pt}
	
	% ─────────────────────────────────────────────────────────────
	% PUNTO 4: PICTOGRAMA Y CIRCULAR
	% ─────────────────────────────────────────────────────────────
	{\large\bfseries\color{valred} 4. Pictograma y Gráfica Circular (25\%)}\\[4pt]
	
	\begin{minipage}[t]{0.48\linewidth}
		\textbf{A) Pictograma:} Venta de Autos. \\
		\textit{Cada \tikz[baseline=-2pt]{\fill[valblue] (0,0) rectangle (0.3,0.3);} vale \textbf{10 autos}. Medio vale \textbf{5}.}
		
		\vspace{4pt}
		\renewcommand{\arraystretch}{1.5}
		\begin{tabular}{|l|l|}
			\hline
			\rowcolor{valblue!15} \textbf{Mes} & \textbf{Autos vendidos} \\ \hline
			Enero & \tikz[baseline=-2pt]{\fill[valblue] (0,0) rectangle (0.3,0.3); \fill[valblue] (0.4,0) rectangle (0.7,0.3);} \\ \hline
			Febrero & \tikz[baseline=-2pt]{\fill[valblue] (0,0) rectangle (0.3,0.3); \fill[valblue] (0.4,0) rectangle (0.7,0.3); \fill[valblue] (0.8,0) rectangle (1.1,0.3); \fill[valblue] (1.2,0) rectangle (1.35,0.3);} \\ \hline
		\end{tabular}
		\vspace{6pt}
		
		\begin{enumerate}[label=\arabic*., itemsep=0.2cm, leftmargin=*]
			\small
			\item ¿Cuántos autos se vendieron en \textbf{Enero}? \hfill\boxfill
			\item ¿Cuántos autos se vendieron en \textbf{Febrero}? \hfill\boxfill
		\end{enumerate}
	\end{minipage}%
	\hfill
	\begin{minipage}[t]{0.48\linewidth}
		\textbf{B) Gráfica Circular:} 200 personas encuestadas sobre su sabor de helado favorito.
		
		\vspace{4pt}
		\begin{minipage}{0.4\linewidth}
			\begin{tikzpicture}[scale=0.8]
				% Vainilla 50%
				\draw[fill=vallightblue, draw=white] (0,0) -- (90:2) arc (90:270:2) -- cycle;
				% Fresa 25%
				\draw[fill=valred!60, draw=white] (0,0) -- (270:2) arc (270:360:2) -- cycle;
				% Chocolate 25%
				\draw[fill=valdark!60, draw=white] (0,0) -- (0:2) arc (0:90:2) -- cycle;
				\node[font=\scriptsize\bfseries, valdark] at (180:1) {50\%};
				\node[font=\scriptsize\bfseries, white] at (-45:1) {25\%};
				\node[font=\scriptsize\bfseries, white] at (45:1) {25\%};
			\end{tikzpicture}
		\end{minipage}%
		\begin{minipage}{0.55\linewidth}
			\begin{enumerate}[start=3, label=\arabic*., itemsep=0.2cm, leftmargin=*]
				\small
				\item Personas que prefieren el helado de \textbf{50\%} (la mitad de 200): \hfill\boxfill[1cm]
				\item Personas que prefieren el de \textbf{25\%}: \hfill\boxfill[1cm]
			\end{enumerate}
		\end{minipage}
	\end{minipage}
	
	\vfill
	\begin{center}
		\textit{\color{valdark!70} ``El éxito es la suma de pequeños esfuerzos repetidos día tras día.''}
	\end{center}
	
\end{document}	